\begin{proof}[Proof of \cref{obj:lem:marked-cubic-mse}]
Fix \(n\ge256\), \(P\in\mathcal M_n\), and \(A\in\{Y,D\}\). All expectations below are under the equal-size source--target sampling law of \(P\). Use the marked-density vector \(F\), the clipped pilot \(\widehat F\), and the corrections \(\mathsf L_A,\mathsf Q_A,\mathsf C_A\) from \cref{obj:def:marked-density-vector,obj:def:cubic-estimator}. Define the training sigma-field and four envelope constants by
\[
\begin{aligned}
\Gamma
&:=\sigma\!\left(
O_r^{\mathrm S},X_r^{\mathrm T}:
r\in\mathcal B_0^{\mathrm S}
=\mathcal B_0^{\mathrm T}
\right),\\
U_{\mathrm{env}}&:=2(1+C_f),&
H_{\mathrm{env}}&:=3(1+C_f)L,\\
V_{\mathrm{bin}}&:=1+C_f+C_f^2,&
M_{\Phi}&:=48(1+C_f)^2(4/c_f)^5.
\end{aligned}
\]
The sampling conditions in
\cref{obj:ass:source-iid,obj:ass:target-iid,obj:ass:sample-independence}
make the three evaluation blocks independent of the training records and mutually independent.

\begin{enumerate}
\item \textit{Envelope and resolution bounds.}
The density, overlap, and arm-mean bounds imply that \(F(x)\) belongs to the clipping rectangle of \cref{obj:lem:observable-marked-reduction} for every \(x\in[0,1]\). The pilot belongs to the same rectangle by \cref{obj:synth_8}. Consequently,
\[
0\le F_i(x),\widehat F_i(x)\le C_f,
\qquad
|F_i(x)-\widehat F_i(x)|\le2C_f
\quad(i\in\mathcal I,\ x\in[0,1]).
\]
Moreover, \cref{obj:lem:observable-marked-reduction} supplies the functional representation and derivative bounds
\[
T_A=\int_0^1\Phi_A(F(x))\,dx,
\qquad
\bigl|\partial_{i_1\cdots i_s}\Phi_A(\widehat F(x))\bigr|
\le M_{\Phi},
\quad 0\le s\le4.
\]
The product seminorm inequality, applied first to the arm densities and then to their marked versions, gives
\[
[q_z]_\beta\le(1+C_f)L,\qquad
[r_{Az}]_\beta
\le C_fL+(1+C_f)L
\le H_{\mathrm{env}},
\]
while \([f_{\mathrm T}]_\beta\le L\le H_{\mathrm{env}}\). Thus
\[
|F_i(x)-F_i(y)|
\le H_{\mathrm{env}}|x-y|^{1/8}
\quad(i\in\mathcal I,\ x,y\in[0,1]).
\]
% lean: CausalSmith.Stat.TransportCaceRoughnuisanceLength.markedDensityVector_holder_modulus

The exact block sizes and the dyadic rules of
\cref{obj:synth_3,obj:synth_4,obj:synth_5,obj:synth_6}
give
\[
m_b\ge\lfloor n/4\rfloor\ge n/4-1\ge n/5>0
\quad(b=0,1,2,3),
\]
and
\[
\frac12m_0^{4/5}\le K_0\le m_0^{4/5},\qquad
\frac12m_0^{4/3}\le K_2\le m_0^{4/3},\qquad
\frac12m_0\le K_3\le m_0.
\]
Indeed, rounding a nonnegative logarithmic exponent down changes its power of two by a factor between \(1/2\) and \(1\). In particular,
\[
K_2\le n^{4/3},\qquad K_3\le n.
\]
The nesting property in \cref{obj:def:cubic-estimator} ensures that the pilot is constant on every correction cell, including its assigned endpoints.
% lean: CausalSmith.Stat.TransportCaceRoughnuisanceLength.block_size_lower
% lean: CausalSmith.Stat.TransportCaceRoughnuisanceLength.pilotResolution_power_bounds
% lean: CausalSmith.Stat.TransportCaceRoughnuisanceLength.quadraticResolution_power_bounds
% lean: CausalSmith.Stat.TransportCaceRoughnuisanceLength.cubicResolution_le_blockSize

\item \textit{Exact projected means.}
For every positive integer \(K\), use the cells \(E_{\ell,K}\) of \cref{obj:synth_7}, with \(0\le\ell<K\), and define the cell-averaging operator on integrable functions by
\[
(\Pi_K\varphi)(x)
:=
K\sum_{\ell=0}^{K-1}\mathbf1\{x\in E_{\ell,K}\}
\int_{E_{\ell,K}}\varphi(y)\,dy,
\qquad x\in[0,1].
\]
For \(K_0\mid K\), define
\[
\begin{aligned}
\delta_i(x)&:=F_i(x)-\widehat F_i(x),\\
\delta^{\mathrm{av}}_{i,K}(x)
&:=(\Pi_KF_i)(x)-\widehat F_i(x),\\
\delta^{\mathrm{av}}_{i,K,\ell}
&:=\delta^{\mathrm{av}}_{i,K}(x_{\ell+1,K}),\\
\eta^{(b)}_{i,K,\ell}
&:=H_{i,K}^{(b)}(x_{\ell+1,K})
-(\Pi_KF_i)(x_{\ell+1,K}),
\qquad b=1,2,3.
\end{aligned}
\]
Here and below the spatial error functions have domain \([0,1]\). The cell-average shifts satisfy
\[
|\delta^{\mathrm{av}}_{i,K,\ell}|
\le2C_f\le U_{\mathrm{env}}.
\]

For \(0\le r<n\), define the observable bin scores and their means by
\[
\mathsf B_{i,K,\ell,r}
:=
W_{i,r}\mathbf1\{X_r^{g(i)}\in E_{\ell,K}\},
\qquad
\pi_{i,K,\ell}:=\int_{E_{\ell,K}}F_i(x)\,dx.
\]
The marked moment identities in \cref{obj:lem:observable-marked-reduction}, together with the boundedness in \cref{obj:ass:outcome-consistency}, give
\[
0\le\mathsf B_{i,K,\ell,r}\le1,\qquad
\mathbb E_P[\mathsf B_{i,K,\ell,r}]=\pi_{i,K,\ell},
\qquad
0\le\pi_{i,K,\ell}\le C_f/K.
\]
Thus
\[
\mathbb E_P[H_{i,K}^{(b)}(x_{\ell+1,K})\mid\Gamma]
=(\Pi_KF_i)(x_{\ell+1,K}),
\qquad
\mathbb E_P[\eta^{(b)}_{i,K,\ell}\mid\Gamma]=0.
\]

Define the three projected corrections and their total by
\[
\begin{aligned}
\overline{\mathsf L}_A
&:=\sum_{i\in\mathcal I}\int_0^1
\partial_i\Phi_A(\widehat F(x))\delta_i(x)\,dx,\\
\overline{\mathsf Q}_A
&:=\frac1{2K_2}
\sum_{i,j\in\mathcal I}\sum_{\ell=0}^{K_2-1}
\partial_{ij}\Phi_A(\widehat F(x_{\ell+1,K_2}))
\delta^{\mathrm{av}}_{i,K_2,\ell}
\delta^{\mathrm{av}}_{j,K_2,\ell},\\
\overline{\mathsf C}_A
&:=\frac1{6K_3}
\sum_{i,j,k\in\mathcal I}\sum_{\ell=0}^{K_3-1}
\partial_{ijk}\Phi_A(\widehat F(x_{\ell+1,K_3}))
\delta^{\mathrm{av}}_{i,K_3,\ell}
\delta^{\mathrm{av}}_{j,K_3,\ell}
\delta^{\mathrm{av}}_{k,K_3,\ell},\\
\overline T_A
&:=\int_0^1\Phi_A(\widehat F(x))\,dx
+\overline{\mathsf L}_A+\overline{\mathsf Q}_A
+\overline{\mathsf C}_A.
\end{aligned}
\]
These are training-measurable quantities. Since the coefficient and pilot are constant on each pilot cell, the linear spatial term is exactly
\[
\overline{\mathsf L}_A
=
\frac1{K_0}\sum_{i\in\mathcal I}\sum_{\ell=0}^{K_0-1}
\partial_i\Phi_A(\widehat F(x_{\ell+1,K_0}))
\delta^{\mathrm{av}}_{i,K_0,\ell}.
\]
The sample covariates belong to \([0,1]\) almost surely, so partitioning the empirical linear term into these cells yields
\[
\mathsf L_A-\overline{\mathsf L}_A
=
\frac1{K_0}\sum_{i\in\mathcal I}\sum_{\ell=0}^{K_0-1}
\partial_i\Phi_A(\widehat F(x_{\ell+1,K_0}))
\eta^{(1)}_{i,K_0,\ell}.
\]

At a correction-cell midpoint,
\[
R^{(b)}_{i,K}
=\delta^{\mathrm{av}}_{i,K,\ell}
+\eta^{(b)}_{i,K,\ell}.
\]
Expanding the quadratic product gives its constant product and three terms containing centered factors; expanding the cubic product gives its constant product and seven such terms. Every nonconstant term has conditional mean zero, because its centered factors belong to distinct evaluation blocks. This remains true when coordinate indices coincide. Therefore, almost surely,
\[
\mathbb E_P[\mathsf L_A\mid\Gamma]=\overline{\mathsf L}_A,\qquad
\mathbb E_P[\mathsf Q_A\mid\Gamma]=\overline{\mathsf Q}_A,\qquad
\mathbb E_P[\mathsf C_A\mid\Gamma]=\overline{\mathsf C}_A,
\]
and hence
\[
\widehat T_A-\overline T_A
=
(\mathsf L_A-\overline{\mathsf L}_A)
+(\mathsf Q_A-\overline{\mathsf Q}_A)
+(\mathsf C_A-\overline{\mathsf C}_A).
\]
The bounded coefficients and finite cell sums supply the integrability needed for these conditional identities.
% lean: CausalSmith.Stat.TransportCaceRoughnuisanceLength.linearProjectionMean_eq_cell_sum
% lean: CausalSmith.Stat.TransportCaceRoughnuisanceLength.linearTerm_sub_projectionMean_eq_centered_cell_sum
% lean: CausalSmith.Stat.TransportCaceRoughnuisanceLength.condExp_linearTerm_eq_linearProjectionMean
% lean: CausalSmith.Stat.TransportCaceRoughnuisanceLength.condExp_quadraticTerm_eq_quadraticProjectionMean
% lean: CausalSmith.Stat.TransportCaceRoughnuisanceLength.condExp_cubicTerm_eq_cubicProjectionMean
% lean: CausalSmith.Stat.TransportCaceRoughnuisanceLength.cubicEstimator_sub_projectedEstimatorMean_eq_centered_corrections

\item \textit{Linear conditional second moment.}
For each coordinate define
\[
\gamma^{\mathrm{lin}}_{i,\ell}
:=\partial_i\Phi_A(\widehat F(x_{\ell+1,K_0})),
\qquad
\mathsf Z^{\mathrm{lin}}_i
:=\frac1{K_0}\sum_{\ell=0}^{K_0-1}
\gamma^{\mathrm{lin}}_{i,\ell}\eta^{(1)}_{i,K_0,\ell}.
\]
Conditional on training, the signed covariance sum for this cell average is the variance of a single weighted record, divided by \(m_1\):
\[
\begin{aligned}
\mathbb E_P[(\mathsf Z^{\mathrm{lin}}_i)^2\mid\Gamma]
&=
\frac1{K_0^2}\sum_{\ell,r=0}^{K_0-1}
\gamma^{\mathrm{lin}}_{i,\ell}\gamma^{\mathrm{lin}}_{i,r}
\operatorname{Cov}_P\!\left(
H_{i,K_0}^{(1)}(x_{\ell+1,K_0}),
H_{i,K_0}^{(1)}(x_{r+1,K_0})
\right)\\
&=
\frac1{m_1}
\operatorname{Var}_P\!\left(
\sum_{\ell=0}^{K_0-1}
\gamma^{\mathrm{lin}}_{i,\ell}
\mathsf B_{i,K_0,\ell,r_0}
\,\middle|\,\Gamma
\right)
\le\frac{M_\Phi^2}{m_1}
\le5M_\Phi^2n^{-1},
\end{aligned}
\]
where \(r_0\) is any record index in evaluation block \(1\), in channel \(g(i)\). The variance inequality follows because the cells are disjoint, each mark lies in \([0,1]\), and therefore the weighted single-record score has absolute value at most \(M_\Phi\). This calculation applies separately to the source and target channels.

Since \(\mathsf L_A-\overline{\mathsf L}_A=\sum_i\mathsf Z_i^{\mathrm{lin}}\), the finite-sum square inequality gives
\[
\begin{aligned}
\mathbb E_P[(\mathsf L_A-\overline{\mathsf L}_A)^2\mid\Gamma]
&\le7\sum_{i\in\mathcal I}
\mathbb E_P[(\mathsf Z_i^{\mathrm{lin}})^2\mid\Gamma]\\
&\le5\cdot7^2M_\Phi^2n^{-1}
\le10\cdot7^2M_\Phi^2n^{-1}.
\end{aligned}
\]
The same finite-sum argument proves square integrability of the centered linear correction.
% lean: CausalSmith.Stat.TransportCaceRoughnuisanceLength.normalized_flatBlockHeight_covariance_sum_le
% lean: CausalSmith.Stat.TransportCaceRoughnuisanceLength.condExp_sq_bounded_linear_cellAverage_le
% lean: CausalSmith.Stat.TransportCaceRoughnuisanceLength.linear_conditional_variance_with_memLp

\item \textit{Quadratic and cubic conditional second moments.}
For two coordinates using the same sample channel, direct expansion of the empirical bin sums gives
\[
\begin{aligned}
&\operatorname{Cov}_P\!\left(
H_{i,K}^{(b)}(x_{\ell+1,K}),
H_{j,K}^{(b)}(x_{r+1,K})
\right)\\
&\qquad=
\frac{K^2}{m_b}
\left\{
\mathbb E_P[
\mathsf B_{i,K,\ell,r_0}\mathsf B_{j,K,r,r_0}]
-\pi_{i,K,\ell}\pi_{j,K,r}
\right\}.
\end{aligned}
\]
For different channels this covariance is zero. On the same cell, bounded marks bound the product moment by \(C_f/K\); on distinct cells the product is zero. It follows that
\[
\left|
\operatorname{Cov}_P\!\left(
H_{i,K}^{(b)}(x_{\ell+1,K}),
H_{j,K}^{(b)}(x_{r+1,K})
\right)
\right|
\le
\begin{cases}
(C_fK+C_f^2)/m_b\le V_{\mathrm{bin}}K/m_b,
&\ell=r,\\
C_f^2/m_b\le V_{\mathrm{bin}}/m_b,
&\ell\ne r.
\end{cases}
\]

To apply these bounds, let \(s\in\{2,3\}\), let
\(\boldsymbol i=(i_1,\ldots,i_s)\in\mathcal I^s\), and let
\(\varnothing\ne\mathcal J\subseteq\{1,\ldots,s\}\). Define
\[
\begin{aligned}
\partial_{\boldsymbol i}\Phi_A
&:=\partial_{i_1\cdots i_s}\Phi_A,\\
\gamma^{(s)}_{\boldsymbol i,\mathcal J,\ell}
&:=
\frac1{s!}
\partial_{\boldsymbol i}\Phi_A(\widehat F(x_{\ell+1,K_s}))
\prod_{j\notin\mathcal J}
\delta^{\mathrm{av}}_{i_j,K_s,\ell},\\
\mathsf Z^{(s)}_{\boldsymbol i,\mathcal J}
&:=
\frac1{K_s}\sum_{\ell=0}^{K_s-1}
\gamma^{(s)}_{\boldsymbol i,\mathcal J,\ell}
\prod_{j\in\mathcal J}\eta^{(j)}_{i_j,K_s,\ell},\\
C_{\mathrm{coef},s}
&:=\frac{M_\Phi(1+U_{\mathrm{env}})^{s-1}}{s!}.
\end{aligned}
\]
Here \(K_s\) means \(K_2\) when \(s=2\) and \(K_3\) when \(s=3\), and an empty product equals \(1\). Since at most \(s-1\) shifts occur,
\[
|\gamma^{(s)}_{\boldsymbol i,\mathcal J,\ell}|
\le C_{\mathrm{coef},s}.
\]
The product expansion from the preceding step now reads
\[
\mathsf Q_A-\overline{\mathsf Q}_A
=\sum_{\boldsymbol i\in\mathcal I^2}
\sum_{\varnothing\ne\mathcal J\subseteq\{1,2\}}
\mathsf Z^{(2)}_{\boldsymbol i,\mathcal J},
\qquad
\mathsf C_A-\overline{\mathsf C}_A
=\sum_{\boldsymbol i\in\mathcal I^3}
\sum_{\varnothing\ne\mathcal J\subseteq\{1,2,3\}}
\mathsf Z^{(3)}_{\boldsymbol i,\mathcal J}.
\]

Put \(\nu:=|\mathcal J|\). Independence of the distinct evaluation blocks gives
\[
\mathbb E_P\!\left[
\prod_{j\in\mathcal J}
\eta^{(j)}_{i_j,K_s,\ell}\eta^{(j)}_{i_j,K_s,r}
\,\middle|\,\Gamma
\right]
=
\prod_{j\in\mathcal J}
\operatorname{Cov}_P\!\left(
H_{i_j,K_s}^{(j)}(x_{\ell+1,K_s}),
H_{i_j,K_s}^{(j)}(x_{r+1,K_s})
\right).
\]
There are \(K_s\) diagonal cell pairs and \(K_s(K_s-1)\) off-diagonal pairs. Expanding the square and applying the covariance bounds therefore yields
\[
\mathbb E_P\!\left[
(\mathsf Z^{(s)}_{\boldsymbol i,\mathcal J})^2
\,\middle|\,\Gamma
\right]
\le
C_{\mathrm{coef},s}^2V_{\mathrm{bin}}^\nu
\frac{5^\nu(K_s^{\nu-1}+1)}{n^\nu}.
\]
This also gives square integrability of every summand.

The constants satisfy \(U_{\mathrm{env}}>4\) and \(V_{\mathrm{bin}}>0\), so
\[
\begin{aligned}
\max\{10V_{\mathrm{bin}},25V_{\mathrm{bin}}^2\}
&\le4(1+U_{\mathrm{env}})^2(1+V_{\mathrm{bin}})^2,\\
\max\{10V_{\mathrm{bin}},25V_{\mathrm{bin}}^2,
125V_{\mathrm{bin}}^3\}
&\le36(1+U_{\mathrm{env}})^2(1+V_{\mathrm{bin}})^3.
\end{aligned}
\]
For \(\nu=1\), the preceding bound has factor \(10V_{\mathrm{bin}}/n\); for \(\nu=2\), it has factor \(25V_{\mathrm{bin}}^2(K_s+1)/n^2\); and for \(\nu=3\), it has factor \(125V_{\mathrm{bin}}^3(K_s^2+1)/n^3\).
Using \(n\ge1\), these inequalities give, for each summand,
\[
\begin{aligned}
\mathbb E_P[(\mathsf Z^{(2)}_{\boldsymbol i,\mathcal J})^2\mid\Gamma]
&\le
M_\Phi^2(1+U_{\mathrm{env}})^4(1+V_{\mathrm{bin}})^2
(n^{-1}+K_2n^{-2}),\\
\mathbb E_P[(\mathsf Z^{(3)}_{\boldsymbol i,\mathcal J})^2\mid\Gamma]
&\le
M_\Phi^2(1+U_{\mathrm{env}})^6(1+V_{\mathrm{bin}})^3
(n^{-1}+K_3n^{-2}+K_3^2n^{-3}).
\end{aligned}
\]
For example, \(K_s+1\le K_s+n\) absorbs the pair contribution, and \(K_3^2+1\le K_3^2+n^2\) absorbs the triple contribution.

There are exactly \(7^s(2^s-1)\) summands: \(7^s\) coordinate tuples, each with \(2^s-1\) nonempty subsets of evaluation positions. Define
\[
\mathcal V_s
:=
7^{2s}(2^s-1)^2M_\Phi^2
(1+U_{\mathrm{env}})^{2s}(1+V_{\mathrm{bin}})^s,
\qquad s=2,3.
\]
Applying the finite-sum square inequality to the displayed expansions gives
\[
\begin{aligned}
\mathbb E_P[(\mathsf Q_A-\overline{\mathsf Q}_A)^2\mid\Gamma]
&\le\mathcal V_2(n^{-1}+K_2n^{-2}),\\
\mathbb E_P[(\mathsf C_A-\overline{\mathsf C}_A)^2\mid\Gamma]
&\le\mathcal V_3(n^{-1}+K_3n^{-2}+K_3^2n^{-3}).
\end{aligned}
\]
% lean: CausalSmith.Stat.TransportCaceRoughnuisanceLength.conditional_variance_centeredSingle_cellAverage
% lean: CausalSmith.Stat.TransportCaceRoughnuisanceLength.conditional_variance_centeredPair_cellAverage
% lean: CausalSmith.Stat.TransportCaceRoughnuisanceLength.conditional_variance_centeredCubic_cellAverage
% lean: CausalSmith.Stat.TransportCaceRoughnuisanceLength.quadratic_conditional_variance_with_memLp
% lean: CausalSmith.Stat.TransportCaceRoughnuisanceLength.cubic_conditional_variance_with_memLp

\item \textit{Fluctuation bound at the common rate.}
The resolution bounds imply
\[
n^{-1}\le n^{-2/3},\qquad
n^{-1}+K_2n^{-2}\le2n^{-2/3},
\]
and
\[
K_3n^{-2}\le n^{-1},\qquad
K_3^2n^{-3}\le n^{-1},\qquad
n^{-1}+K_3n^{-2}+K_3^2n^{-3}\le3n^{-2/3}.
\]
Enlarging the nonnegative quadratic and cubic constants gives
\[
\begin{aligned}
\mathbb E_P[(\mathsf L_A-\overline{\mathsf L}_A)^2\mid\Gamma]
&\le10\cdot7^2M_\Phi^2n^{-2/3},\\
\mathbb E_P[(\mathsf Q_A-\overline{\mathsf Q}_A)^2\mid\Gamma]
&\le60\mathcal V_2n^{-2/3},\\
\mathbb E_P[(\mathsf C_A-\overline{\mathsf C}_A)^2\mid\Gamma]
&\le310\mathcal V_3n^{-2/3}.
\end{aligned}
\]
Define
\[
\mathcal V_{\mathrm{tot}}
:=3\{10\cdot7^2M_\Phi^2+60\mathcal V_2+310\mathcal V_3\}.
\]
Applying the three-term square inequality to the exact fluctuation identity, and then taking expectations, gives
\[
\mathbb E_P[(\widehat T_A-\overline T_A)^2\mid\Gamma]
\le\mathcal V_{\mathrm{tot}}n^{-2/3},
\qquad
\mathbb E_P[(\widehat T_A-\overline T_A)^2]
\le\mathcal V_{\mathrm{tot}}n^{-2/3}.
\]
The component square-integrability statements also give
\[
\widehat T_A-\overline T_A\in L^2(P).
\]
% lean: CausalSmith.Stat.TransportCaceRoughnuisanceLength.quadratic_conditional_variance_sample_rate
% lean: CausalSmith.Stat.TransportCaceRoughnuisanceLength.cubic_conditional_variance_sample_rate
% lean: CausalSmith.Stat.TransportCaceRoughnuisanceLength.linear_conditional_variance_sample_rate
% lean: CausalSmith.Stat.TransportCaceRoughnuisanceLength.centered_corrections_conditional_second_moment_sample_rate
% lean: CausalSmith.Stat.TransportCaceRoughnuisanceLength.centered_corrections_expected_conditional_second_moment_sample_rate
% lean: CausalSmith.Stat.TransportCaceRoughnuisanceLength.cubicEstimator_fluctuation_second_moment_sample_rate

\item \textit{Pilot moments.}
For a fixed coordinate and cell, define the centered training-record scores and their variance by
\[
\xi_{i,K,\ell,r}
:=\mathsf B_{i,K,\ell,r}-\pi_{i,K,\ell},
\qquad
\sigma_{i,K,\ell}^2
:=\mathbb E_P[\xi_{i,K,\ell,r}^2].
\]
Bounded bin scores give
\[
|\xi_{i,K,\ell,r}|\le1,\qquad
\mathbb E_P[\xi_{i,K,\ell,r}]=0,\qquad
\sigma_{i,K,\ell}^2
\le C_f/K+(C_f/K)^2
\le V_{\mathrm{bin}}/K.
\]
The centered training histogram is exactly
\[
H_{i,K}^{(0)}(x_{\ell+1,K})
-(\Pi_KF_i)(x_{\ell+1,K})
=
\frac K{m_0}\sum_{r\in\mathcal B_0^{g(i)}}
\xi_{i,K,\ell,r}.
\]
Independence therefore gives its second-moment bound
\[
\mathbb E_P\!\left[
\left\{H_{i,K}^{(0)}(x_{\ell+1,K})
-(\Pi_KF_i)(x_{\ell+1,K})\right\}^2
\right]
\le V_{\mathrm{bin}}K/m_0.
\]

For the eighth moment, expand the eighth power of the centered record sum and group terms by their equality pattern among the eight record positions. A pattern containing a singleton has expectation zero. Every surviving pattern has between one and four blocks. For a pattern with \(j\) blocks, boundedness of the centered scores bounds each block moment of order at least two by \(\sigma_{i,K,\ell}^2\), and there are at most \(m_0^j\) choices of its record indices. Its absolute contribution is consequently at most
\[
(m_0\sigma_{i,K,\ell}^2)^j
\le m_0\sigma_{i,K,\ell}^2
+(m_0\sigma_{i,K,\ell}^2)^4,
\qquad 1\le j\le4.
\]
There are at most \(8!\) equality patterns: encoding each block as a cycle in increasing order embeds the set of partitions into the permutations of eight positions. Thus
\[
\mathbb E_P\!\left[
\left(\sum_{r\in\mathcal B_0^{g(i)}}\xi_{i,K,\ell,r}\right)^8
\right]
\le8!\left\{
(m_0\sigma_{i,K,\ell}^2)^4
+m_0\sigma_{i,K,\ell}^2
\right\}.
\]
Scaling by \((K/m_0)^8\) yields
\[
\mathbb E_P\!\left[
\left\{H_{i,K}^{(0)}(x_{\ell+1,K})
-(\Pi_KF_i)(x_{\ell+1,K})\right\}^8
\right]
\le8!\left\{
V_{\mathrm{bin}}^4(K/m_0)^4
+V_{\mathrm{bin}}(K/m_0)^7
\right\}.
\]
% lean: CausalSmith.Stat.TransportCaceRoughnuisanceLength.iidBlockHeight_eighth_moment_le

Averaging the Hölder modulus over a cell gives
\[
|(\Pi_KF_i)(x)-F_i(x)|
\le H_{\mathrm{env}}K^{-1/8},
\qquad x\in E_{\ell,K}.
\]
Each clipping interval contains \(F_i(x)\), so clipping contracts the distance to that value. Splitting the histogram error at its cell mean, using
\((a+b)^2\le2(a^2+b^2)\) and
\((a+b)^8\le2^7(a^8+b^8)\) for nonnegative \(a,b\), gives
\[
\begin{aligned}
\mathbb E_P[\delta_i(x)^2]
&\le2V_{\mathrm{bin}}K_0/m_0
+2H_{\mathrm{env}}^2K_0^{-1/4},\\
\mathbb E_P[|\delta_i(x)|^8]
&\le2^7\,8!\left\{
V_{\mathrm{bin}}^4(K_0/m_0)^4
+V_{\mathrm{bin}}(K_0/m_0)^7
\right\}
+2^7H_{\mathrm{env}}^8K_0^{-1}.
\end{aligned}
\]
The dyadic bounds imply
\[
\begin{aligned}
K_0/m_0&\le5^{1/5}n^{-1/5},&
K_0^{-1/4}&\le2^{1/4}5^{1/5}n^{-1/5},\\
(K_0/m_0)^4&\le5^{4/5}n^{-4/5},&
(K_0/m_0)^7&\le5^{7/5}n^{-7/5}
\le5^{7/5}n^{-4/5},\\
K_0^{-1}&\le2\cdot5^{4/5}n^{-4/5}.
\end{aligned}
\]
Define
\[
\begin{aligned}
C_{\mathrm{pilot},2}
&:=2V_{\mathrm{bin}}5^{1/5}
+2^{5/4}H_{\mathrm{env}}^25^{1/5},\\
C_{\mathrm{pilot},8}
&:=2^7\,8!\left\{
V_{\mathrm{bin}}^45^{4/5}
+V_{\mathrm{bin}}5^{7/5}\right\}
+2^8H_{\mathrm{env}}^85^{4/5}.
\end{aligned}
\]
We have obtained, for every \(x\in[0,1]\) and \(i\in\mathcal I\),
\[
\mathbb E_P[\delta_i(x)^2]
\le C_{\mathrm{pilot},2}n^{-1/5},
\qquad
\mathbb E_P[|\delta_i(x)|^8]
\le C_{\mathrm{pilot},8}n^{-4/5}.
\]
Finally, define the spatial pilot errors by
\[
\mathsf E_{i,\mathrm{pilot}}:=\int_0^1|\delta_i(x)|\,dx.
\]
Cauchy--Schwarz on the unit interval and integration of the pointwise bounds give
\[
\begin{aligned}
\mathbb E_P[\mathsf E_{i,\mathrm{pilot}}^2]
&\le\mathbb E_P\int_0^1\delta_i(x)^2\,dx
\le C_{\mathrm{pilot},2}n^{-1/5},\\
\mathbb E_P\int_0^1|\delta_i(x)|^8\,dx
&\le C_{\mathrm{pilot},8}n^{-4/5}.
\end{aligned}
\]
Joint measurability and the bound \(|\delta_i|\le2C_f\) justify these integrations.
% lean: CausalSmith.Stat.TransportCaceRoughnuisanceLength.pilot_second_moment
% lean: CausalSmith.Stat.TransportCaceRoughnuisanceLength.pilot_eighth_moment
% lean: CausalSmith.Stat.TransportCaceRoughnuisanceLength.pilot_error_L1_second_moment

\item \textit{Spatial Taylor remainder.}
For \(s=1,2,3\), define the spatial Taylor terms by
\[
\mathsf T^{\mathrm{Taylor}}_{A,s}
:=
\frac1{s!}
\sum_{(i_1,\ldots,i_s)\in\mathcal I^s}
\int_0^1
\partial_{i_1\cdots i_s}\Phi_A(\widehat F(x))
\prod_{j=1}^s\delta_{i_j}(x)\,dx.
\]
In particular,
\(\mathsf T^{\mathrm{Taylor}}_{A,1}
=\overline{\mathsf L}_A\).

The remainder estimate follows from an exact ratio identity. For \(z=0,1\), define
\[
\Psi_{Az}(\mathbf w)
:=\frac{w_t w_{r_{Az}}}{w_{q_z}},
\qquad
\mathbf w\in\mathbb R^7,\quad w_{q_z}>0,
\]
using the named-coordinate convention of \cref{obj:synth_14}. For each \(x\in[0,1]\), direct expansion and reduction to a common denominator give
\[
\begin{aligned}
&\Psi_{Az}(F(x))-\Psi_{Az}(\widehat F(x))
-\sum_{s=1}^3\frac1{s!}
D^s\Psi_{Az}(\widehat F(x))
[\delta(x),\ldots,\delta(x)]\\
&\quad=
\frac{
\delta_{q_z}(x)^2
\{\widehat F_{q_z}(x)\delta_t(x)
-\widehat F_t(x)\delta_{q_z}(x)\}
\{\widehat F_{q_z}(x)\delta_{r_{Az}}(x)
-\widehat F_{r_{Az}}(x)\delta_{q_z}(x)\}
}{
\widehat F_{q_z}(x)^4 F_{q_z}(x)
}.
\end{aligned}
\]
Here
\[
\delta(x):=(\delta_0(x),\ldots,\delta_6(x)).
\]
Both denominator coordinates are at least \(c_f/4\). Each numerator factor in braces is bounded by \(C_f\) times the sum of the absolute increments in its two coordinates. Bounding these sums and \(|\delta_{q_z}|\) by the full coordinate sum therefore gives
\[
\left|
\Psi_{Az}(F)-\Psi_{Az}(\widehat F)
-\sum_{s=1}^3\frac1{s!}
D^s\Psi_{Az}(\widehat F)[\delta,\ldots,\delta]
\right|
\le C_f^2(4/c_f)^5
\left(\sum_{i\in\mathcal I}|\delta_i|\right)^4.
\]
Since \(\Phi_A=\Psi_{A1}-\Psi_{A0}\), define its pointwise remainder by
\[
\begin{aligned}
\mathcal R_{A,4}(x)
:=\Phi_A(F(x))-\Phi_A(\widehat F(x))
-\sum_{s=1}^3\frac1{s!}
\sum_{(i_1,\ldots,i_s)\in\mathcal I^s}
\partial_{i_1\cdots i_s}\Phi_A(\widehat F(x))
\prod_{j=1}^s\delta_{i_j}(x).
\end{aligned}
\]
The two ratio bounds and \(C_f^2\le(1+C_f)^2\) imply
\[
|\mathcal R_{A,4}(x)|
\le2(1+C_f)^2(4/c_f)^5
\left(\sum_i|\delta_i(x)|\right)^4
=
\frac{M_\Phi}{24}
\left(\sum_i|\delta_i(x)|\right)^4.
\]
% lean: CausalSmith.Stat.TransportCaceRoughnuisanceLength.cubicRatio_cubic_taylor_remainder_eq
% lean: CausalSmith.Stat.TransportCaceRoughnuisanceLength.Phi_cubic_taylor_remainder_abs_le

Define the integrated remainder and its majorant by
\[
\mathsf R_{\mathrm{rem},A}
:=\int_0^1\mathcal R_{A,4}(x)\,dx,
\qquad
\mathsf G_{\mathrm{rem}}
:=\frac{M_\Phi}{24}
\int_0^1\left(\sum_i|\delta_i(x)|\right)^4dx.
\]
Then
\[
|\mathsf R_{\mathrm{rem},A}|\le\mathsf G_{\mathrm{rem}},
\qquad
0\le\mathsf G_{\mathrm{rem}}
\le\frac{M_\Phi}{24}(14C_f)^4.
\]
Cauchy--Schwarz on \([0,1]\), followed by the finite-sum eighth-power inequality, gives
\[
\begin{aligned}
\mathbb E_P[\mathsf G_{\mathrm{rem}}^2]
&\le
(M_\Phi/24)^2
\mathbb E_P\int_0^1
\left(\sum_i|\delta_i(x)|\right)^8dx\\
&\le
(M_\Phi/24)^2\,7^7
\sum_{i\in\mathcal I}\mathbb E_P\int_0^1|\delta_i(x)|^8dx\\
&\le
(M_\Phi/24)^2\,7^8C_{\mathrm{pilot},8}n^{-4/5}.
\end{aligned}
\]
Define
\[
\mathcal B_{\mathrm{rem}}
:=(M_\Phi/24)^2\,7^8C_{\mathrm{pilot},8}.
\]
Since \(n^{-4/5}\le n^{-2/3}\),
\[
\mathbb E_P[\mathsf G_{\mathrm{rem}}^2]
\le\mathcal B_{\mathrm{rem}}n^{-2/3}.
\]
Integrating the defining Taylor identity also yields
\[
T_A=
\int_0^1\Phi_A(\widehat F(x))\,dx
+\sum_{s=1}^3\mathsf T^{\mathrm{Taylor}}_{A,s}
+\mathsf R_{\mathrm{rem},A}.
\]
% lean: CausalSmith.Stat.TransportCaceRoughnuisanceLength.pilot_remainder_integral_abs_le_majorant
% lean: CausalSmith.Stat.TransportCaceRoughnuisanceLength.pilotRemainderMajorant_integrable_sq_and_moment
% lean: CausalSmith.Stat.TransportCaceRoughnuisanceLength.pilotRemainderMajorant_sample_rate
% lean: CausalSmith.Stat.TransportCaceRoughnuisanceLength.spatialTaylorRemainder_eq_transport_sub_coordinate_terms

\item \textit{Exact spatial projection errors.}
For \(K_0\mid K\), define
\[
\varepsilon_{i,K}(x):=F_i(x)-(\Pi_KF_i)(x),
\qquad x\in[0,1].
\]
On each cell,
\[
\delta_i=\delta^{\mathrm{av}}_{i,K}+\varepsilon_{i,K},
\qquad
\int_{E_{\ell,K}}\varepsilon_{i,K}(x)\,dx=0,
\qquad
|\varepsilon_{i,K}(x)|\le H_{\mathrm{env}}K^{-1/8}.
\]
The pilot and all averaged shifts are constant on that cell. Hence any cell product containing exactly one averaging residual has integral zero. More explicitly, for any function \(\vartheta:\mathbb R^7\to\mathbb R\) and any finite list of coordinates \(i_1,\ldots,i_s\),
\[
\int_{E_{\ell,K}}
\vartheta(\widehat F(x))
\left\{\prod_{j=1}^s\delta^{\mathrm{av}}_{i_j,K}(x)\right\}
\varepsilon_{i,K}(x)\,dx=0.
\]
This includes \(s=0\), with the empty product equal to \(1\).
% lean: CausalSmith.Stat.TransportCaceRoughnuisanceLength.spatial_pilot_single_residual_cancellation

Define the two projection errors by
\[
\mathsf R_{\mathrm{quad},A}
:=\mathsf T^{\mathrm{Taylor}}_{A,2}
-\overline{\mathsf Q}_A,
\qquad
\mathsf R_{\mathrm{cub},A}
:=\mathsf T^{\mathrm{Taylor}}_{A,3}
-\overline{\mathsf C}_A.
\]
The nested-cell identities identify these as differences between the true spatial Taylor terms and the exact midpoint projection means. For the quadratic error, expanding the two factors and cancelling the single-residual terms gives
\[
\mathsf R_{\mathrm{quad},A}
=
\frac12\sum_{i,j\in\mathcal I}\sum_{\ell=0}^{K_2-1}
\partial_{ij}\Phi_A(\widehat F(x_{\ell+1,K_2}))
\int_{E_{\ell,K_2}}
\varepsilon_{i,K_2}(x)\varepsilon_{j,K_2}(x)\,dx.
\]
Thus
\[
|\mathsf R_{\mathrm{quad},A}|
\le\frac{7^2}{2}M_\Phi H_{\mathrm{env}}^2K_2^{-1/4}.
\]
Using the lower dyadic bound,
\[
K_2^{-1/2}
\le2^{1/2}5^{2/3}n^{-2/3}.
\]
Define
\[
\mathcal B_{\mathrm{quad}}
:=7^4M_\Phi^2H_{\mathrm{env}}^4
2^{1/2}5^{2/3}.
\]
Squaring the absolute bound and enlarging its factor \(1/4\) gives the pointwise inequality
\[
\mathsf R_{\mathrm{quad},A}^2
\le\frac{7^4M_\Phi^2H_{\mathrm{env}}^4}{4}K_2^{-1/2}
\le\mathcal B_{\mathrm{quad}}n^{-2/3}.
\]
% lean: CausalSmith.Stat.TransportCaceRoughnuisanceLength.quadraticCellProjectionBias_eq_spatial_sub_projectionMean
% lean: CausalSmith.Stat.TransportCaceRoughnuisanceLength.quadraticCellProjectionBias_abs_le
% lean: CausalSmith.Stat.TransportCaceRoughnuisanceLength.quadraticCellProjectionBias_sq_sample_rate

For the cubic error, work on a \(K_3\)-cell. Expanding the three factors leaves, after single-residual cancellation,
\[
\begin{aligned}
&\int_{E_{\ell,K_3}}
\left\{
\delta_i\delta_j\delta_k
-\delta^{\mathrm{av}}_{i,K_3,\ell}
 \delta^{\mathrm{av}}_{j,K_3,\ell}
 \delta^{\mathrm{av}}_{k,K_3,\ell}
\right\}dx\\
&\quad=
\delta^{\mathrm{av}}_{i,K_3,\ell}
\int_{E_{\ell,K_3}}\varepsilon_{j,K_3}\varepsilon_{k,K_3}\,dx
+\delta^{\mathrm{av}}_{j,K_3,\ell}
\int_{E_{\ell,K_3}}\varepsilon_{i,K_3}\varepsilon_{k,K_3}\,dx\\
&\qquad+
\delta^{\mathrm{av}}_{k,K_3,\ell}
\int_{E_{\ell,K_3}}\varepsilon_{i,K_3}\varepsilon_{j,K_3}\,dx
+\int_{E_{\ell,K_3}}
\varepsilon_{i,K_3}\varepsilon_{j,K_3}\varepsilon_{k,K_3}\,dx.
\end{aligned}
\]
To retain the pilot error, use the exact cell-average identity
\[
\frac{|\delta^{\mathrm{av}}_{i,K_3,\ell}|}{K_3}
=
\left|\int_{E_{\ell,K_3}}\delta_i(x)\,dx\right|
\le\int_{E_{\ell,K_3}}|\delta_i(x)|\,dx.
\]
Consequently, the absolute value of the preceding cell integral is at most
\[
H_{\mathrm{env}}^2K_3^{-1/4}
\int_{E_{\ell,K_3}}(|\delta_i|+|\delta_j|+|\delta_k|)\,dx
+\frac{H_{\mathrm{env}}^3K_3^{-3/8}}{K_3}.
\]
Multiplying by the cubic coefficient, bounded by \(M_\Phi/6\), and summing over cells and coordinate triples gives
\[
|\mathsf R_{\mathrm{cub},A}|
\le
C_{\mathrm{bias},2}K_3^{-1/4}
\sum_{i\in\mathcal I}\mathsf E_{i,\mathrm{pilot}}
+C_{\mathrm{bias},3}K_3^{-3/8},
\]
where
\[
C_{\mathrm{bias},2}
:=\frac{M_\Phi}{2}7^2H_{\mathrm{env}}^2,
\qquad
C_{\mathrm{bias},3}
:=\frac{M_\Phi}{6}7^3H_{\mathrm{env}}^3.
\]
Define the square majorant by
\[
\mathsf S_{\mathrm{cub}}
:=
2C_{\mathrm{bias},2}^2K_3^{-1/2}
\,7\sum_{i\in\mathcal I}\mathsf E_{i,\mathrm{pilot}}^2
+2C_{\mathrm{bias},3}^2K_3^{-3/4}.
\]
The two-term square inequality and the coordinate-sum square inequality imply
\[
\mathsf R_{\mathrm{cub},A}^2\le\mathsf S_{\mathrm{cub}}.
\]
The pilot \(L^1\) moment bound therefore gives
\[
\mathbb E_P[\mathsf S_{\mathrm{cub}}]
\le
2C_{\mathrm{bias},2}^2\,7^2C_{\mathrm{pilot},2}
K_3^{-1/2}n^{-1/5}
+2C_{\mathrm{bias},3}^2K_3^{-3/4}.
\]
The lower bound \(K_3\ge n/10\) yields
\[
\begin{aligned}
K_3^{-1/2}n^{-1/5}
&\le2^{1/2}5^{1/2}n^{-7/10}
\le2^{1/2}5^{1/2}n^{-2/3},\\
K_3^{-3/4}
&\le2^{3/4}5^{3/4}n^{-3/4}
\le2^{3/4}5^{3/4}n^{-2/3}.
\end{aligned}
\]
Define
\[
\mathcal B_{\mathrm{cub}}
:=
2\left\{
C_{\mathrm{bias},2}^2\,7^2C_{\mathrm{pilot},2}
2^{1/2}5^{1/2}
+C_{\mathrm{bias},3}^2\,2^{3/4}5^{3/4}
\right\}.
\]
We conclude that
\[
\mathbb E_P[\mathsf S_{\mathrm{cub}}]
\le\mathcal B_{\mathrm{cub}}n^{-2/3}.
\]
% lean: CausalSmith.Stat.TransportCaceRoughnuisanceLength.cubic_midpoint_cell_projection_L1_bound
% lean: CausalSmith.Stat.TransportCaceRoughnuisanceLength.cubic_spatial_projection_bias_abs_le_L1
% lean: CausalSmith.Stat.TransportCaceRoughnuisanceLength.cubic_bias_resolution_terms_sample_rate

\item \textit{Squared bias of the projected mean.}
The integrated Taylor identity and the definitions of the projection errors give the exact signed identity
\[
\overline T_A-T_A
=
-\left(
\mathsf R_{\mathrm{rem},A}
+\mathsf R_{\mathrm{quad},A}
+\mathsf R_{\mathrm{cub},A}
\right).
\]
Using the three-term square inequality and the preceding majorants gives
\[
(\overline T_A-T_A)^2
\le
3\left\{
\mathsf G_{\mathrm{rem}}^2
+\mathcal B_{\mathrm{quad}}n^{-2/3}
+\mathsf S_{\mathrm{cub}}
\right\}.
\]
The right-hand side is integrable, by the remainder bound and the spatial pilot moment bounds. Define
\[
\mathcal B_{\mathrm{tot}}
:=3\left(
\mathcal B_{\mathrm{rem}}
+\mathcal B_{\mathrm{quad}}
+\mathcal B_{\mathrm{cub}}
\right).
\]
Integration yields
\[
\mathbb E_P[(\overline T_A-T_A)^2]
\le\mathcal B_{\mathrm{tot}}n^{-2/3},
\qquad
\overline T_A-T_A\in L^2(P).
\]
% lean: CausalSmith.Stat.TransportCaceRoughnuisanceLength.combined_spatial_projection_bias_second_moment
% lean: CausalSmith.Stat.TransportCaceRoughnuisanceLength.projectedEstimatorMean_sub_transport_eq_spatial_bias
% lean: CausalSmith.Stat.TransportCaceRoughnuisanceLength.projectedEstimatorMean_squared_bias_sample_rate

\item \textit{Orthogonality and final assembly.}
Define the fluctuation and projected bias by
\[
\mathsf R_{\mathrm{fluc},A}:=\widehat T_A-\overline T_A,
\qquad
\mathsf R_{\mathrm{bias},A}:=\overline T_A-T_A.
\]
The exact conditional means imply
\[
\mathbb E_P[\mathsf R_{\mathrm{fluc},A}\mid\Gamma]=0.
\]
The projected bias is training-measurable, and both quantities are square-integrable. Their product is therefore integrable, and pulling out the training-measurable factor gives
\[
\begin{aligned}
\mathbb E_P[
\mathsf R_{\mathrm{bias},A}\mathsf R_{\mathrm{fluc},A}]
&=
\mathbb E_P\!\left[
\mathsf R_{\mathrm{bias},A}
\mathbb E_P[\mathsf R_{\mathrm{fluc},A}\mid\Gamma]
\right]
=0.
\end{aligned}
\]
Expanding the square consequently gives the exact risk decomposition
\[
\begin{aligned}
\mathbb E_P[(\widehat T_A-T_A)^2]
&=
\mathbb E_P[\mathsf R_{\mathrm{fluc},A}^2]
+\mathbb E_P[\mathsf R_{\mathrm{bias},A}^2]\\
&\le
(\mathcal V_{\mathrm{tot}}+\mathcal B_{\mathrm{tot}})n^{-2/3}.
\end{aligned}
\]
Finally, substitution of the defined constants gives
\[
\begin{aligned}
\mathcal V_{\mathrm{tot}}+\mathcal B_{\mathrm{tot}}
&=
3\left(
\mathcal B_{\mathrm{rem}}
+\mathcal B_{\mathrm{quad}}
+\mathcal B_{\mathrm{cub}}
\right)
+3\left\{
10\cdot7^2M_\Phi^2
+60\mathcal V_2+310\mathcal V_3
\right\}\\
&=C_{\mathrm{mse}},
\end{aligned}
\]
with exactly the expanded expression in the statement. Each term is finite and explicitly computable from \((c_f,C_f,L)\), since \(c_f>0\). All bounds apply to both response choices and every \(P\in\mathcal M_n\), proving the asserted uniform inequality.
% lean: CausalSmith.Stat.TransportCaceRoughnuisanceLength.condExp_cubicEstimator_sub_projectedEstimatorMean_eq_zero
% lean: CausalSmith.Stat.TransportCaceRoughnuisanceLength.integral_projected_bias_mul_fluctuation_eq_zero
% lean: CausalSmith.Stat.TransportCaceRoughnuisanceLength.cubicEstimator_mse_eq_fluctuation_add_projected_bias
% lean: CausalSmith.Stat.TransportCaceRoughnuisanceLength.cubicEstimator_mse_sample_rate
% lean: CausalSmith.Stat.TransportCaceRoughnuisanceLength.marked_cubic_mse
\end{enumerate}
\end{proof}
